

\chapter{Review of Related Literature}\label{sec:RRL}
The evolution of financial market prediction using computer vision is slowly becoming a popular tool used in trading. Previous studies surrounding market prediction can be categorized into three categories. To start, we have numerical models that are solely trained on technical indicators. Another approach is by using computer vision, which analyzes chart-based or image representations of price series. Lastly, multimodal fusion models that cover signals such as prices, technical indicators, and macroeconomic variables to improve prediction performance.

Chakole et al.\ \cite{chakole2021convolutional} introduced a Deep Trend Following (DTF) strategy that uses CNN prediction with market trend signals. The approach focuses on generating trading decisions based on the current and predicted future price of stocks. Results showed that their model provided a much better performance than traditional buy-and-hold models. Similar to this method, Saud et al.\ \cite{saud2024technical} used technical indicators such as MACD, DMI, and KST with LSTM models and GRU networks that can be used to develop intelligent strategies to predict stock trading signals. Their study showed that the MACD-based approach was the safest and most effective approach among other models. These results reinforce the role of technical indicators as a strong foundation for quantitative trading models.

Computer vision offers another effective approach for financial market prediction. 
Unlike traditional numerical models, image-based methods convert temporal market information into structured visual representations such as candlestick images, bar-chart images, Gramian Angular Fields, recurrence plots, or technical indicator grids. 
This transformation is useful because it allows convolutional filters to process local spatial patterns that correspond to temporal relationships in the original series. 
Wang and Oates showed that time-series data can be encoded as images using Gramian Angular Fields and Markov Transition Fields, allowing CNN-based models to learn discriminative patterns from visualized temporal structures \cite{wang2015encoding}. 
Hatami et al. also showed that recurrence plots can transform time-series data into two-dimensional texture-like images, making time-series classification suitable for CNN-based image recognition \cite{hatami2018classification}. 

In financial applications, this idea is consistent with the way traders visually interpret chart patterns before making decisions. 
Cohen et al. \cite{cohen2020trading} treated trading as an image classification task by transforming financial time-series data into visual forms. 
Sezer and Ozbayoglu \cite{sezer2018algorithmic} further demonstrated that financial indicators can be converted into two-dimensional image-like representations for CNN-based algorithmic trading. 
Chen and Tsai \cite{chen2020encoding} used Gramian Angular Fields to encode candlestick data as images and showed that CNNs can recognize candlestick patterns from these transformed representations. 
Liu and Kang \cite{liu2023automated} further supported this direction by showing that candlestick-based representations can be used in cryptocurrency trading models and can outperform purely numeric representations in their experimental setting.
These studies support the use of image-based encoding because temporal signals become organized into spatial patterns that CNN filters can analyze.

The effectiveness of translating financial data into images is strongly demonstrated by a progression of studies from Sezer \cite{sezer2018algorithmic}. Building on earlier work that converted indicators into images (CNN-TA) and used bar-chart images (CNN-BI), Maratkhan et al.\ \cite{maratkhan2021deep} presented a highly optimized pipeline. Their architectures (CNN--TIC, EO--CNN--TIC, ResNet--TA) clustered technical indicators into multi-channel images and used an optimized trading layer to achieve remarkable results: ``Buy/Sell'' recall scores up to $0.94/0.96$ and average annual returns on DOW-30 up to $36.70\%$, substantially outperforming buy-and-hold and their own reproduced baselines. In addition, Wu \cite{wu2021empowering} specifically highlighted that computer vision techniques can address technical analysis objectives like predicting bullish/bearish trends. Then in 2019, Chen et al.\ \cite{chen2020encoding} developed approaches to automatically recognize candlestick patterns with over $80\%$ accuracy using convolutional neural networks, providing strong evidence for the effectiveness of visual representation in trading.

Studies on financial market forecasting also utilized multimodal fusion frameworks to enhance prediction accuracy by integrating different data sources. Wang et al.\ \cite{wang2019aggregating} pioneered a model-independent framework that successfully incorporated trading volume, intraday return series, and investors’ social media emotions to forecast market trends. This study demonstrates that combining multiple information sources significantly outperforms single-method approaches in financial market prediction. Ma et al.\ \cite{ma2022macroeconomic} proved that macroeconomic attention indices (MAI) can predict stock market returns by components that are extracted using partial least squares (PLS) and the LASSO method. Chauhan et al.\ \cite{chauhan2025exploring} had a similar study that used an ARDL model to reveal how macroeconomic indicators like Foreign Institutional Investment (FII) and Economic Policy Uncertainty (EPU) influence stock returns over both short- and long-term market behavior.

Finally, researchers have fused textual and numerical data. Mo et al.\ \cite{mo2024mof} proposed a system that addresses key limitations in traditional stock prediction models by integrating policy and stock review texts with price data. The mechanism uses MacBERT to generate feature vectors and combine them with stock price data to be processed on a neural network. Similarly, Rahul Modak et al.\ \cite{Modak2023MarketSA} developed a multimodal transformer architecture that integrates earnings call transcripts, social media content, and technical market indicators, achieving $87.2\%$ sentiment classification accuracy compared to $78.4\%$ for unimodal baselines.

